\section{Results}

\subsection{How often the record cannot support its own verdict}
\label{sec:results-waterbase}
Waterbase records a below-quantification flag and the limit of quantification
in force for each station, substance and
year\src{scripts/22\_waterbase\_external.py}. Keeping river water only, since
an inland-surface-water standard does not apply to lakes or coastal waters,
leaves \num{4190833} station-years from \num{37} reporting countries.

\paragraph{Censoring is the normal case, not an edge case.}
\textbf{\num{44.9}\% of all samples lie below the limit of quantification.}
Figure~\ref{fig:europe-map} maps the stations: the pattern follows national
borders, not river basins.

\begin{figure}[htbp]\centering
\includegraphics[width=0.86\textwidth]{fig02_europe_map.pdf}
\caption{The \num{7501} monitoring stations in the graph that carry
coordinates, in the equal-area projection the EEA publishes European water data
in (ETRS89-LAEA), coloured by the share of their reporting authority's
station-years that record no quantification limit. The value is a property of the
\emph{authority}, not of the station --- every station in a country carries the
same one --- so each national cluster is labelled with its code and its rate
rather than left to the colour scale alone. Countries are drawn lowest rate
first, so no reporter is hidden beneath another. No coastline is shown: no
boundary file is distributed with this work, and the outline of Europe here is
the monitoring network itself. The pattern follows national borders rather than
river basins, which places the defect in reporting practice rather than in
chemistry or hydrology. Every plotted station, with its country, is in
\texttt{supplementary/figure\_data/fig02\_europe\_map.csv}. Generated by
\texttt{scripts/90\_figures.py}.}
\label{fig:europe-map}
\end{figure}

\paragraph{The bound is missing for a quarter of the record.}
Where censoring is declared the limit is almost always recorded. The failure is
elsewhere: \textbf{\num{25.5}\,\% of station-years declare nothing at all}, no
flag and no limit, so a zero and a measured trace are indistinguishable and no
substitution rule can be chosen. These are \texttt{UnresolvedObservation} at
continental scale.

\begin{figure}[htbp]\centering
\includegraphics[width=0.86\textwidth]{fig03_country_practice.pdf}
\caption{Each point is a reporting country with at least \num{5000}
station-years; area is proportional to that count. The relationship runs the
wrong way for a chemical explanation: authorities recording least (right) are
found in breach least (bottom), because a limit that was never reported cannot
be tested against the standard. This is not a correlation between two
independent quantities and is not offered as one --- the horizontal axis governs
the \emph{denominator} of the vertical one, so the downward slope is the
signature of that dependence rather than a finding about chemistry. What it
establishes is one thing: the measured failure rate is a lower bound, and it is
loosest exactly where reporting is worst. Generated by \texttt{scripts/90\_figures.py}.}
\label{fig:country-practice}
\end{figure}

\paragraph{More than a third of European monitoring fails the legal
performance criterion.}
Article~4(1) of Directive 2009/90/EC requires a limit of quantification at or
below \num{30}\,\% of the relevant environmental quality standard before a
method's results may be used~\cite{eu2009qaqc}. \textbf{\num{40.1}\% of station-years assessed under a European standard were produced by a method that
does not meet it.}

\paragraph{Article~3(3b) is routinely unenforceable.}
Of the station-years for which a European standard exists, \textbf{\num{17.5}\% report a below-quantification result whose limit exceeds that standard}, which the Directive
requires to be set aside. For three pyrethroid insecticides --- deltamethrin,
bifenthrin and esfenvalerate --- the figure is \num{100}\%: none has ever, in
this record, been measured by a method capable of deciding its own standard.
Figure~\ref{fig:loqeqs}a ranks substances by the stricter Article~4(1)
criterion, on which a fourth --- clothianidin, a neonicotinoid --- also reaches
\num{100}\,\%.

\paragraph{The failure scales with the standard, not with the substance.}
The substances in Figure~\ref{fig:loqeqs}a share no chemistry, only the
\emph{magnitude of the limit written for them}. By decade of the
annual-average standard, the share whose quantification limit exceeds it rises
from \num{0.6}\,\% at order \num{10}~\si{\micro\gram\per\litre} to
\num{81.8}\,\% at \num{e-4}~\si{\micro\gram\per\litre}, over decades
carrying between \num{5} and \num{16} substances each
(Figure~\ref{fig:loqeqs}b). The two decades below that reach \num{100}\,\%
but rest on four substances and on one, so the trend is read from the decades
that can carry it.
Directive~(EU) 2026/805 moves the analyte list into exactly those decades:
\num{8} of the \num{24} substances it adds with an annual-average standard
put it below \num{e-3}~\si{\micro\gram\per\litre}
(\num{33}\,\%), against \num{7} of \num{35} on the legacy list
(\num{20}\,\%). The defect therefore grows unless analytical capability moves
down with the threshold, and nothing in the record suggests it has.

\begin{figure}[htbp]\centering
\includegraphics[width=0.86\textwidth]{fig04_loq_vs_eqs.pdf}
\caption{\textbf{(a)} The \emph{fourteen} substances with the highest share of
station-years whose quantification limit exceeds \num{30}\,\% of the
environmental quality standard, among those with at least \num{1000} assessable
station-years and at least one such station-year. The substances are selected
\emph{by} the quantity plotted, so the shaded band is the \num{95}\,\% Wilson
interval: ranking puts the least certain estimates on show, and the band is what
separates a substance assessed \num{25960} times from one assessed \num{1098}
times. Every substance meeting the criterion is in
\texttt{supplementary/figure\_data/fig04\_loq\_vs\_eqs.csv} with its interval,
and a column recording which were plotted. \textbf{(b)} The same failure against
the \emph{magnitude} of the standard rather than the identity of the substance,
over every decade holding at least \num{500} assessable station-years: the lower
the limit the law writes, the less often the monitoring can decide it. Between
\num{5} and \num{16} substances stand behind each decade from
\num{e1} down to \num{e-4}~\si{\micro\gram\per\litre}; the two lowest are
thin, and the circled point is a decade holding \emph{one} substance, whose rate
is that substance's rate and not a property of the decade. The composition of
every decade ships beside the figure. Generated by
\texttt{scripts/90\_figures.py}.}
\label{fig:loqeqs}
\end{figure}

\subsection{The record is not one era, and the three failures do not move together}
\label{sec:results-era}
Every row carries its own reference year, so the audit can be read year by
year\src{scripts/22\_waterbase\_external.py}. \num{48.8}\,\% of the river
record is dated \num{2015} or later, and the three failures behave differently
over time --- but not in the same way, and not gradually
(Figure~\ref{fig:by-year}).

\textbf{One failure was repaired on a date.} The share of station-years
declaring neither a censoring flag nor a limit falls from \num{31.8}\,\% in
\num{2012} to \num{0.0}\,\% in \num{2013}, and stays there for the
\num{12} years to \num{2024}: not a trend, a step. \textbf{The other two did
not move.} The share of assessments from a method failing the \num{30}\,\%
criterion is \num{46.3}\,\% in \num{2006} and \num{47.7}\,\% in
\num{2024}; the share whose limit exceeds the standard outright, \num{28.1}\,\%
and \num{28.2}\,\%. Substitution at source never moves at all: censored rows
carrying a positive value stay between \num{97.2}\,\% and \num{100}\,\%
across every year the record can speak about. Twelve years after the
record-keeping defect was repaired, the analytical defect is where it was.

\paragraph{The recent record is the stronger case, not the weaker one.} Before
\num{2013} an objection was available: the failures might be artefacts of
incomplete reporting. After it that objection is gone --- every station-year
carries its limit --- and the analytical failure is unchanged. Read on
\num{2024} alone, the last year in the record, \num{47.7}\,\% of assessments
still come from a method failing the performance criterion and \num{28.2}\,\%
from one whose limit exceeds the standard outright. What the paper reports is
therefore not a defect of an archive.

\begin{figure}[htbp]\centering
\includegraphics[width=0.86\textwidth]{fig08_by_year.pdf}
\caption{\textbf{(a)} The three failures by reference year, over the
\num{29} years holding at least \num{20000} station-years. The dashed line is
read from the data, not chosen: it marks the first year after which no reporter
omits the limit again. \textbf{(b)} The denominator, without which panel (a)
misleads: before \num{2005} almost no row carries a limit, so the legal criteria
read \num{0}\,\% for want of anything to test rather than because the
monitoring met them. Every year, plotted or not, is in
\texttt{supplementary/figure\_data/fig08\_by\_year.csv}. Generated by
\texttt{scripts/90\_figures.py}.}
\label{fig:by-year}
\end{figure}

Splitting at \num{2015} rather than at the break keeps the comparison
conservative and lets it be read against the era table
(Table~\ref{tab:era}).

\begin{table}[htbp]\centering
\caption{The audit split at \num{2015}. Shares of station-years, except the
legal criteria, which are over station-years for which a European standard
exists, and substitution, which is over censored rows. Full table in
\texttt{supplementary/S10\_by\_era.csv}, and by reporting country in
\texttt{S11\_era\_by\_country.csv}.}
\label{tab:era}
\begin{tabular}{lrr}
\toprule
 & before \num{2015} & \num{2015} onwards \\
 & (\num{2147102}) & (\num{2043731}) \\
\midrule
neither a censoring flag nor a limit & \num{49.7}\,\% & \num{0.0}\,\% \\
quantification limit above the standard & \num{22.7}\,\% & \num{21.8}\,\% \\
fails the \num{30}\,\% performance criterion & \num{38.5}\,\% & \num{41.4}\,\% \\
censored rows carrying a positive value & \num{99.6}\,\% & \num{99.3}\,\% \\
\bottomrule
\end{tabular}
\end{table}

\textbf{The record-keeping failure is solved}: station-years declaring neither
a flag nor a limit fall from \num{49.7}\,\% to \num{0.0}\,\%. \textbf{The
analytical failure does not move.} Its earlier figure is lower only because so
little of the old record can be tested: where no limit was recorded, no method
can be found wanting, so the rate is a lower bound.
\textbf{Substitution at source has not moved at all}: \num{99.3}\,\% of
censored rows still carry a positive number in place of the non-detection.

\paragraph{The fall is practice, not composition.} The set of reporters differs
between the periods, so a falling rate proves nothing on its own. The \num{8}
authorities reporting in both go to \num{0.0}\,\%, Denmark from \num{88.1}\,\%
(\texttt{supplementary/S11\_era\_by\_country.csv}), with no change to the
standard they report under.

\paragraph{What this does not license.} Those \num{8} are the only authorities
with a post-\num{2015} river record here, and \num{29} of the \num{37}
reporters have none, so any statement about \emph{current European practice}
rests on a handful of programmes. What holds without qualification is narrower:
where the record-keeping defect was repaired, the reporter repaired it, and the
analytical defect survived the repair.

\subsection{Two modelling choices the record decides}
\label{sec:results-modelling}
The record decides two commitments from Section~\ref{sec:methods-vocab} that
separate CENSO from the vocabularies in
Section~\ref{sec:results-gap}\src{scripts/22\_waterbase\_external.py}: where
the quantification limit lives, and what a standard written over a sum means.

\paragraph{The limit is not a property of the instrument.}
If it were, the same substance at the same station would carry the same limit
each year. Of the \num{557956} station--substance pairs reported in more than one
year, \textbf{\num{56.4}\,\% report more than one limit}. Whether the method
changed with it cannot be read from the record: only \num{13654} of those pairs
--- \num{2.4}\,\% --- report an analytical method code at all, and \num{3356}
of those that do report more than one. In \textbf{\num{21043} pairs the change
crosses the standard}: the same station and substance is decidable in one year
and not the next. A vocabulary that binds the limit to a sensor cannot express
that; one that binds it to the run at least has somewhere to put the provenance
this record is missing.

\paragraph{A sum standard is not a limit on any of its members.}
Annex~I states four limits over a group rather than a substance --- a
component-based standard resting on concentration
addition~\cite{backhaus2012mixtures,posthuma2019mixtures} --- and a sum from an
incompletely measured basket is an underestimate the arithmetic never reveals;
CENSO carries \texttt{requiresCompleteGroup} for it. All four resolve from the
Annex, and the completeness rate falls with the size of the basket: heptachlor
and its epoxide (\num{2} members) are both measured in \num{69.6}\,\% of the
station-years that touch them, the cyclodiene pesticides (\num{4}) in
\num{88.0}\,\%, the polyaromatic hydrocarbons (\num{9}) in \num{36.0}\,\%
--- and the sum over the pesticides list, which resolves to \num{75} members,
in \textbf{\num{0.0}\,\%}: not one of the \num{26846} station-years that
touch that group measures all of it. A standard defined over seventy-five
substances that no station-year ever measures completely cannot be assessed by
anyone, and a threshold column would compare the partial sum against it and
return compliance.

\subsection{The ontology applied to open data}
\label{sec:results-wbabox}
Expressing a Waterbase subset in CENSO\src{scripts/23\_waterbase\_abox.py}
tests whether its distinctions survive a schema written by someone else: every
class is assigned from EEA columns.

\paragraph{The assessment is the population, not a sample.}
The decision is arithmetic on four reported fields, applied to all
\num{696168} station-years whose substance carries a European
standard\src{scripts/22\_waterbase\_external.py}. \textbf{\num{43.8}\,\% are not
decidable from the record as reported} --- one outcome,
\texttt{IndeterminateCompliance}, reached for four distinct reasons the
vocabulary keeps apart as its subclasses: \num{18.8}\,\%
\texttt{PreconditionUnmet}, where the standard is defined on a quantity the
record does not report (Section~\ref{sec:results-precondition});
\num{17.5}\,\% \texttt{MethodInsufficient}, where a below-quantification result
carries a limit above the standard, which Article~3(3b) requires to be set
aside; \num{6.5}\,\% \texttt{BoundNotEstablished}, with no bound to compare at
all; \num{1.0}\,\% \texttt{PossibleExceedance}, where the interval straddles the
limit. A two-valued schema has nowhere to put any of them and collapses all four
into compliance.

The taxonomy is not free ornament, and the cheaper alternative shows why: a
schema with a single indeterminacy flag reports the same \num{43.8}\,\% and
loses the four. They are not degrees of one defect but different objects, and
they call for different repairs --- a schema field for the covariates, a better
laboratory, a reported limit, a reported uncertainty --- of which only one is
addressed by anything an analyst can do to the data afterwards.

\paragraph{One reason comes from the standard rather than from the record.}
\texttt{PossibleExceedance} requires an interval around a quantified result,
and no European release reports a per-result measurement uncertainty,
aggregated or disaggregated. Article~4(1) caps that uncertainty at
\num{50}\,\% ($k=2$) at the level of the standard, so the widest lawful
interval around a reported value $x$ is $x \pm 0.5\,T$, straddling the
standard when $0.5\,T < x < 1.5\,T$ --- a band around $T$, not an
extrapolation. What is reported is a bound, not an estimate.

\paragraph{What a two-valued pipeline reports for the same rows.}
Entering a non-detection at half its quantification limit is not a convention an
analyst chose. Article~5(1) of Directive 2009/90/EC requires it: ``where the
amounts of physico-chemical or chemical measurands in a given sample are below
the limit of quantification, the measurement results shall be set to half of the
value of the limit of quantification concerned''~\cite{eu2009qaqc}. The same
instrument that sets the performance criterion also mandates the substitution,
and the substitution pushes results \emph{above} the standard: of the \num{171285}
exceedances such a pipeline reports, \num{112034} exist only because the
substitution lifted a censored result past the limit, \num{30140} are asserted
against a standard that is not defined on the quantity the record reports, and
\num{12142} rest on a number whose censoring status the record never states.
\num{14505} are exceedances the law can affirm --- quantified, above the
standard, with an applicable standard and an interval that does not straddle it.
A further \num{2458} are quantified above the standard and set aside all the
same, because the widest interval Article~4(1) permits straddles it. Article~3(3b)
sets aside no quantified result: it reaches a value reported as below the
quantification limit, and a mean quantified above the standard is evidence of
exceedance whatever the method's limit was. Zero, half the limit and the limit
itself report \num{59251}, \num{171285} and \num{200708} exceedances over the
same records --- a \num{3.4}-fold range in the enforcement-relevant number,
selected by a rule the data does not determine
(Figure~\ref{fig:verdicts}).

\paragraph{The graph is a materialisation, not the analysis.}
A reservoir sample of \num{40000} station-years is expressed as RDF
(\num{501062} triples): a pure-Python triple store cannot hold four million
rows, and inflating it would misreport what was reasoned over. \num{11.7}\% of
its observations are \texttt{UnresolvedObservation}, a state no other surveyed
vocabulary can express.

\begin{figure}[htbp]\centering
\includegraphics[width=0.86\textwidth]{fig05_verdicts.pdf}
\caption{The same \num{696168} assessments --- every station-year whose
substance carries a European standard, not a sample --- read four ways. The top
three rows are the identical records under the three substitution conventions
in routine use; they disagree about the number of exceedances by a factor of
\num{3.4}, and \num{14505} are exceedances the law can affirm. The bottom
row reports the outcomes the law distinguishes instead of absorbing them into
compliance; \emph{possible exceedance} is a quantified result whose widest
lawful uncertainty interval straddles the standard. The plotted values are in
\texttt{supplementary/figure\_data/fig05\_verdicts.csv}. Generated by
\texttt{scripts/90\_figures.py}.}
\label{fig:verdicts}
\end{figure}

\paragraph{What the graph adds, and what it does not.}
The graph returns the same method-insufficiency share as the streaming count
--- \num{17.2}\,\% on the sample against \num{17.5}\,\% on the population,
inside the \num{0.19}-point sampling error --- but that agreement is a
consistency check rather than corroboration: both routes call one decision procedure, held in a single module
precisely so the two cannot drift. What the check establishes is that
expressing the record in the vocabulary loses nothing --- every row classified
in RDF lands where the counter puts it. What it does not establish is that the
ontology was needed to count: it was not. The counting could always have been
done. What only the graph does is answer \emph{which} assessments are
unsupportable, under which limit, for which substance, in one query.

\subsection{A standard that cannot be applied to what the record reports}
\label{sec:results-precondition}
Two Annex~I footnotes make a threshold conditional rather than absolute. Footnote
12 states that the standards for lead and nickel ``refer to bioavailable
concentrations of the substances''; footnote 9 that cadmium's ``vary depending on
the hardness of the water as specified in five class categories''. Bioavailable
normalisation needs hardness, dissolved organic carbon and pH; the hardness class
needs hardness. WISE-6 reports none of them on the row.

So for these three substances there is no comparison to make. Not a strict one
and not a lenient one: the number in Annex~I and the number in the record are not
measurements of the same quantity, and asking whether the method's limit reaches
\num{30}\,\% of the standard is asking about the wrong quantity. The released
package attaches the two conditions to the six thresholds they govern, and the
assessment returns \texttt{PreconditionUnmet} ahead of every other test ---
before the value, before the limit, before Article~3(3b).

\textbf{That is \num{131068} assessments, \num{18.8}\,\% of every station-year
a European standard reaches}\src{scripts/22\_waterbase\_external.py}. They are
not a marginal stratum: cadmium, lead and nickel are among the most reported
substances in the record. A two-valued pipeline decides all of them anyway, and
decides them in both directions --- \num{30140} are reported as exceedances of a
standard that is not defined on what was measured, the remainder as compliance
with it.

This is the commitment a threshold column cannot hold at all. The other three
could in principle be bolted onto a table: a censoring flag, a limit per row, a
verdict enumeration. A precondition is a property of the \emph{standard} --- it
travels with the regulation, not with the measurement --- so a column of
thresholds has nowhere to record that a threshold was inapplicable, and no way to
distinguish that from a threshold that was met. The vocabulary named lead, nickel
and cadmium in its class comments before this record was read; the record returns
\num{131068} assessments where those comments decide the outcome.

This stratum is metals, and only metals, because Annex~I attaches its
applicability conditions there and nowhere else. It is one of the four reasons an
assessment can be undecidable, not the paper's subject: censoring itself is far
worse among the organic micropollutants, where over three-quarters of samples
fall below the limit against under half for the Annex~I metals
(Section~\ref{sec:disc-lipinski}), and where the substances the record can least
often decide --- pyrethroids, neonicotinoids, pharmaceuticals, bisphenol~A ---
are concentrated (Figure~\ref{fig:loqeqs}a). A condition on the analysed
fraction would reach every metal rather than three: Article~3(1) defines the
metal standards on the dissolved phase and WISE-6 reports the analysed matrix,
so such a condition would be checkable in both directions rather than merely
unsatisfiable. The vocabulary does not declare one, because no footnote in the
consolidated Annex~I creates it and a class no package can instantiate is a
promise rather than a term; it is the next thing the machinery should decide.

\subsection{The substitution has already happened before anyone downloads
the data}
\label{sec:results-asreported}
Sections~\ref{sec:results-waterbase} and~\ref{sec:results-wbabox} treat
substitution as something an analyst does to a record. The ontology's SHACL
shapes, run over the record \emph{as reported} with none of our pipeline's
repair\src{scripts/26\_source\_conformance.py}, show that it is not.

\textbf{\num{55.0}\% of river station-years declare a result below the limit of
quantification and, in the same row, state a positive value for it.} A
non-detection permits any concentration down to zero; a positive lower bound is
a number someone chose, and the reporting authority chose it before
publication. Counted over the censored rows alone rather than over the whole
record, the share is \num{99.4}\,\%: where a row is flagged at all, it almost
always carries a substituted number as well.

This changes who the finding is about. Where the flag survives, as in
Waterbase, the substitution is visible and can be undone, which makes the
counterfactual in Section~\ref{sec:results-wbabox} computable. Where it does
not --- \num{25.5}\% of the record, by Section~\ref{sec:results-waterbase} ---
the same substitution cannot be detected, let alone reversed.

\paragraph{The shapes are the deliverable, not the demonstration.}
A deterministic sample of \num{20000} concentration rows validated against the
published shapes yields \textbf{\num{14137} violations}, all of one kind: a
result below the quantification limit carrying a positive number for it. A
published constraint catches the substitution, not a bespoke script.

An earlier version of the shapes required every method to state a limit of
\emph{detection}, so they fired on every row: WISE-6 has no field for one, and
Article~3(3b) of Directive 2008/105/EC and Article~4(1) of Directive 2009/90/EC
are written about the quantification limit. Our own graph satisfied that
constraint by manufacturing the datum at $\mathrm{LOD}=\mathrm{LOQ}/3$. The
shapes now require \emph{a} limit, and the ABox states the one the source
reports. The distinction the old constraint reached for is real and is made in
Section~\ref{sec:methods-noise}: a non-detection establishes
$[0,\mathrm{LOD}]$, the record at best $[0,\mathrm{LOQ}]$, a weaker and
different claim.

\subsection{The same observations under two regulations}
\label{sec:results-dual}
Pluggable regulation packages let one set of observations be assessed under
several jurisdictions without
reimplementation\src{scripts/24\_dual\_regulation.py}, as
Section~\ref{sec:methods-why} claims: both packages are read as data, and only
the file supplying the threshold changes.

Of the \num{1488677} river station-years regulated by at least one instrument,
\num{81.9}\% change outcome when the package is swapped, but most of that is
coverage: one jurisdiction regulating a substance the other does not, a fact
about the legal texts rather than about any measurement. Where \emph{both}
regulate and only the numeric standard differs, \textbf{\num{17.7}\% (95\,\% CI \num{17.6}--\num{17.9}\,\%) still change outcome}, mostly the \num{47999} the
European standard sets aside as method-insufficient and the looser Turkish
limit can decide.

An assessment that is \emph{method-insufficient} under one instrument and
decidable under another is not expressible if the threshold is a column:
method adequacy is a relation between the limit and the standard. Both verdicts
can be held at once, on thresholds with different IRIs, and neither overwrites
the other. Turkish standards are not in force over European rivers; the
comparison is counterfactual by construction and demonstrates that the
machinery is regulation-independent, never that it assesses Türkiye.

\subsection{A standard defined on a different unit of observation}
\label{sec:results-mac}
Annex~I states two standards for most priority substances: an annual average,
defined against an annual mean, and a maximum allowable concentration, defined
against an \emph{individual measurement}. Everything above evaluates the first,
on the aggregated release, where one row \emph{is} an annual mean. The second
could be approached there --- the release carries an annual maximum --- but not
answered: that column is absent for part of the record, carries its own
below-quantification flag for less of it, and gives no per-sample denominator,
so the share of samples a method cannot decide is not recoverable from it. The
vocabulary holds the two standards apart, and conflating them is an
inconsistency.

Reading the disaggregated release --- \num{96597294} individual sample records
--- makes it computable\src{scripts/27\_mac\_exceedance.py}. Of the \num{8148785}
river samples with such a standard and a convertible unit, \num{74702} exceed it outright and \textbf{\num{794418} (\num{9.7}\,\%) are undecidable}: not
detected, with a quantification limit \emph{above} the standard, so no result
the method could have returned would have settled it. Article~3(3b) requires
those to be set aside; a two-valued pipeline records every one as compliant.

\paragraph{The two standards disagree, and a column can hold only one.}
Of the \num{66285} station-years assessable under both,
\num{25.6}\,\% receive different verdicts: in \textbf{\num{1.1}\,\% --- \num{698} station-years --- the published annual mean is compliant while
a sample behind it breaches the maximum allowable concentration}
(Figure~\ref{fig:two-thresholds}a), in the rest the mean exceeds while no
single sample does. Neither verdict is wrong and neither supersedes the other;
Annex~I asks both. A schema storing one threshold per substance cannot record
which it holds.

\begin{figure}[htbp]\centering
\includegraphics[width=\textwidth]{fig07_two_thresholds.pdf}
\caption{What a threshold column cannot hold, on real observations. Each panel
cross-tabulates the same observations assessed twice; grey is agreement, colour
is the verdict the second assessment changes \emph{into}, and the outlined cell
is the one the text names. \textbf{(a)} The same station-year under the two
standards Annex~I states for it: an annual average, on the published mean, and
a maximum allowable concentration, on the samples behind that mean.
\textbf{(b)} The same observation under two jurisdictions' packages; the
Turkish standards are not in force over European rivers, so this demonstrates
the machinery rather than assessing either jurisdiction. Its
\emph{no bound established} row and column are empty except on the diagonal,
and that is the result rather than an artefact: a record that establishes no
bound establishes none under either instrument, so the one outcome that cannot
move is the one caused by the reporting rather than by the law. The plotted counts are
in \texttt{supplementary/figure\_data/fig07\_two\_thresholds.csv}. Generated by
\texttt{scripts/90\_figures.py}.}
\label{fig:two-thresholds}
\end{figure}

The disaggregated release does not hold the samples behind every aggregated
station-year, so these shares describe the stratum readable at both units and not
the whole record: \num{66285} of the \num{376843} station-years that carry an
annual-average verdict and a maximum-allowable standard --- \num{17.6}\,\% ---
have samples behind them in the disaggregated release. Whether that stratum
resembles the rest is not something the release lets us test, since the missing
station-years are missing precisely their samples; the per-substance and
per-country coverage is in
\texttt{supplementary/S12\_mac\_exceedance.csv} so a reader can see which
substances it rests on. A censored sample is decided
by its quantification limit, never by the number the reporter wrote in its
place, the artefact Section~\ref{sec:results-asreported} exposes.

\subsection{Is the between-country spread an artefact?}
\label{sec:results-confounders}
The share of station-years carrying neither a censoring flag nor a limit ranges
from \num{3.3}\% to \num{97.4}\% between countries reporting under the same
directive (Figure~\ref{fig:country-practice}). The countries differ in
substances, years and Member State status, so the objection is
testable\src{scripts/25\_country\_confounders.py}. WISE-6 defines the field
identically for every reporter, and a national standard changes which results
breach a limit, not whether the limit was written down. Three adjustments ---
the pooled European substance mix, the substances most of these countries
report, and EU membership --- each leave the spread intact, with
Kendall's~$\tau$ between the raw and standardised orderings at \num{0.89};
per-country values are in supplementary table~S5. Restricting to recent years,
in Section~\ref{sec:results-era}, is not a fourth test, because too few
countries report after \num{2015}. A spread stable under three independent
corrections is a property of reporting practice.

\subsection{Comparison with published ontologies}
\label{sec:results-gap}
% generated by scripts/07_verify_gap_table.py
% GENERATED by scripts/07_verify_gap_table.py -- do not edit.
\begin{table*}[htbp]\centering
\caption{Concepts present in published water and observation ontologies, decided by parsing each ontology file rather than its documentation. A dash means the file parsed and the concept was absent; \texttt{?} means it could not be retrieved, and no claim is made. The \emph{LOD bound to} column is the discriminating one: a detection limit declared as a sensor capability cannot state that a given measurement fell below it. The two rows marked \emph{this work} are CENSO's own modules and are not among the assessed vocabularies, which is why the text counts eighteen where the table has twenty rows.}
\label{tab:gap}
\scriptsize\setlength{\tabcolsep}{3.5pt}
\begin{tabular}{l r l cccccc}
\toprule
Ontology & Ent. & LOD bound to & det.\ limit & censoring & undecid. & threshold & interval & applic. \\
\midrule
SOSA/SSN & 37 & -- & -- & -- & -- & -- & -- & -- \\
SOSA core & 53 & -- & -- & -- & -- & -- & -- & -- \\
GeoSPARQL & 0 & -- & -- & -- & -- & -- & -- & -- \\
QUDT schema & 239 & -- & -- & -- & -- & \checkmark & -- & \checkmark \\
SAREF core & 131 & -- & -- & -- & -- & \checkmark & -- & -- \\
SAREF4WATR v2.1.1 & 95 & -- & -- & -- & -- & \checkmark & -- & -- \\
SSN-system & 44 & the sensor & \checkmark & -- & -- & -- & -- & -- \\
OWL-Time & 78 & -- & -- & -- & -- & -- & -- & -- \\
WHOW water-monitoring & 57 & -- & -- & -- & -- & -- & \checkmark & -- \\
DoCE (Rio Doce WQ) & 168 & -- & -- & -- & -- & -- & -- & \checkmark \\
InWaterSense core & 189 & -- & -- & -- & -- & \checkmark & \checkmark & -- \\
InWaterSense regulations & 144 & -- & -- & -- & -- & -- & -- & -- \\
InWaterSense pollutants & 9 & -- & -- & -- & -- & -- & -- & -- \\
WAM-ONTO & 921 & -- & -- & -- & -- & -- & -- & -- \\
SewerNet & 145 & -- & -- & -- & -- & -- & -- & -- \\
CHMO (chemical methods) & 3438 & the method & \checkmark & -- & \checkmark & -- & \checkmark & -- \\
AFO (Allotrope) & 3864 & -- & -- & -- & -- & \checkmark & \checkmark & \checkmark \\
STATO (statistics) & 1162 & -- & -- & -- & -- & \checkmark & \checkmark & -- \\
ENVO & 7208 & -- & -- & -- & -- & -- & -- & -- \\
ExO (exposure) & 195 & -- & -- & -- & -- & -- & -- & -- \\
\textbf{CENSO (this work)} & 50 & \textbf{the result} & \checkmark & \checkmark & \checkmark & \checkmark & \checkmark & \checkmark \\
\textbf{CENSO-REG (this work)} & 17 & -- & -- & -- & -- & \checkmark & -- & -- \\
Project ontology 2018 & 123 & unclear & \checkmark & -- & -- & -- & -- & -- \\
\bottomrule
\end{tabular}
\end{table*}

\num{21} ontology files were assessed by parsing each one rather than reading
its documentation\src{scripts/07\_verify\_gap\_table.py}: \num{20} published
vocabularies plus one unpublished domain ontology as a non-standard baseline
(Table~\ref{tab:gap}). The set deliberately reaches outside the water domain.
CENSO claims applicability to any measurement process with a detection limit,
and a domain-independence claim tested inside one domain is not tested, so the
comparison includes the vocabularies built for the laboratory --- CHMO, the
\textsc{obo} Chemical Methods Ontology; AFO, the Allotrope Foundation Ontology,
built by a pharmaceutical consortium for analytical instrument data; and STATO,
the \textsc{obo} statistics ontology, where censoring is a routine concept in
survival analysis.

\textbf{No ontology except CENSO carries a censoring status.} That column is
empty for all \num{21}, and it is the one the argument rests on. It is also the
column the matcher was made hardest to fail: the patterns admit
\emph{censored}, \emph{left-censored}, \emph{right-censored},
\emph{left-truncated}, \emph{non-detect}, \emph{not detected},
\emph{not quantifiable} and \emph{below \dots\ limit} bare, and STATO --- where
censoring is routine in survival analysis --- declares no term for it and does
not mention it in a single definition. One ontology besides CENSO scores for a
not-determinable outcome, and it should not: CHMO's two hits are
\texttt{ambiguous synonym}, an \textsc{obo} annotation property describing the
quality of a synonym, and \texttt{double quantum transitions for finding
unresolved lines}, a technique in which the unresolved things are spectral
lines. The cell is left standing and the discrepancy published beside the table,
because a pattern narrowed until a competitor loses a cell would make the
comparison an assertion again.

\paragraph{The patterns are deliberately generous, and the asymmetry is the
point.} Every widening can only add a \emph{yes} to a competitor's row: CENSO's
own cells are already filled, so no adjustment to the matcher can flatter this
work. The lists were widened after the first version was found to match on our
vocabulary's spellings and miss other people's --- it did not match
\emph{limit of quantitation}, the pharmacopoeial spelling a laboratory ontology
is most likely to use, nor \emph{below detection limit}, the most ordinary
phrasing there is, nor \emph{minimum value} spelled out. A claim that survives a
matcher tuned in the competitors' favour is worth something; one that depends on
spelling is not.

Two columns need reading carefully, because a keyword match is not a modelling
commitment. Two vocabularies score for an interval-valued result, but only WHOW
attaches the interval to the \emph{result}: InWaterSense's range properties
describe the range a \emph{device} can report, and AFO's \texttt{lower bound}
and \texttt{upper bound} are subclasses of \texttt{bound role} --- roles borne
by a participant, with no property placing an interval on a measured datum.
Those cells are left scored, and the discrepancy published beside the table
rather than corrected away: a pattern adjusted until a competitor scores the way
we want would make the table an assertion again.

\paragraph{The sharpest case is CHMO, and it is not a water ontology.}
\textbf{CHMO declares \texttt{limit of detection} (CHMO:0002801) and
\texttt{limit of quantification} (CHMO:0002802), both defined from the IUPAC
Gold Book} --- so the gap is not that chemistry lacks the concept. CHMO binds
them to the \emph{method}, as a \texttt{figure of merit} of an assay:

\begin{center}
\texttt{limit of detection} $\sqsubseteq$ \texttt{figure of merit} $\sqsubseteq$
\texttt{assay output} $\sqsubseteq$ \texttt{data entity}
\end{center}

\noindent That is a better place than SSN-system's, which makes the limit a
\texttt{SystemProperty} of a \emph{sensor}, and it is still not a place from
which any particular result can be called censored. None of CHMO's
object properties relates a limit to a result; they are the \textsc{bfo}/\textsc{ro}
relations (\texttt{has participant}, \texttt{is specified output of},
\texttt{has characteristic}). The limit is known about the assay and forgotten
about the measurement --- the same failure as the sensor case, one step closer.
This is the result we would most have expected to lose, and it is the reason the
claim is worth making: the concept exists, is well defined, is cited to the
IUPAC Gold Book, and still nothing can say that \emph{this} number is a bound.

Profiling for sampling versus sensing terminology separates the vocabularies
cleanly: SAREF core scores \num{0}/\num{22} sampling to sensing terms, CENSO
\num{3}/\num{0}. \textbf{The field is dominated by vocabularies built for
in-situ sensing}, in which a detection limit is properly a device property, and
the few sampling-aware ones carry no limits at all.

Five further ontologies described only in papers (WaWO+, OPO, InWaterSense's
expert system, Ahmedi et al., Acquirium) were assessed from their full
text\src{scripts/11\_assess\_competitor\_papers.py}: in none do \emph{detection
limit}, \emph{LOD}, \emph{LOQ}, \emph{censoring} or \emph{non-detection} occur.
This is weaker evidence than parsing a file and is labelled as such.

\subsection{What a term list cannot show: the vocabularies cannot separate the
readings}
\label{sec:results-separation}
A gap table records which vocabulary declares which term, and invites an
answer we think a referee is entitled to give: a missing term is a missing
convenience, not a missing capability, and anyone may mint a subclass. The
answer to that is logical rather than lexical. A vocabulary distinguishes two
situations only if some sentence it can write is true of one and false of the
other. So we take one real record, write down the readings it admits, and ask
of each vocabulary whether the readings it can \emph{express} still
differ\src{scripts/13\_separation.py}.

The witness is the row of Figure~\ref{fig:geometry}(b): imidacloprid at a
Spanish river station, reported below a quantification limit of
\num{0.01}~\si{\micro\gram\per\litre}, carrying \num{0.01} as its value,
against an annual-average standard of
\num{0.0068}~\si{\micro\gram\per\litre}. Nothing about it is constructed --- the
standard lies inside the censored interval, and the reported value is the limit
itself, which is what substitution at source looks like in this release. Those
four numbers admit three readings: the censoring is honoured and the result is
the interval $[0,\mathrm{LOQ}]$ (\textbf{R1}); the flag is dropped and the
number is a measurement, which a numeric comparison calls an exceedance
(\textbf{R2}); or the number arrives with neither flag nor limit, the
\num{25.5}\,\% case (\textbf{R3}). R1 and R2 are the pair the paper turns on:
the same four numbers, two different claims.

Each reading is translated into each vocabulary, generously: a fact survives if
the vocabulary has \emph{any} term in the concept family it needs, and every
vocabulary is granted an observation and a reported number without being
checked for them. The families are the gap table's own columns, so nothing is
decided here that parsing the files did not already decide.

\begin{quote}
\textbf{Proposition.} R1 and R2 differ only in facts whose families are
\emph{censoring}, \emph{interval-valued result} and \emph{undecidable verdict}.
A vocabulary declaring no term in those families translates both readings to
the same set of facts; hence for every sentence $\varphi$ in that vocabulary,
R1 entails $\varphi$ exactly when R2 does.
\end{quote}

\textbf{\num{20} of the \num{21} vocabularies collapse R1 into R2.} The one
exception is CHMO, and it is an artefact of the generous matcher rather than a
capability: CHMO is credited with a not-determinable outcome on the strength of
\texttt{ambiguous synonym} and \texttt{unresolved lines}, so R1 receives a
\texttt{no\_verdict} fact that R2 does not. Read strictly it collapses with the
rest, and it has no censoring term either way. We report the number the generous
matcher gives. For the other \num{20}, no query, rule or
reasoner built on them returns a different answer for a non-detection than for
a measurement at the same number --- not because none has been written, but
because none can be. SSN-system and the \num{2018} baseline separate R1 from
R3, but only by noticing that a limit was stated somewhere; neither can say
that \emph{this} result fell below it, which is what the discriminating column
of Table~\ref{tab:gap} records. Two vocabularies with interval-valued results,
WHOW and InWaterSense core, would separate R1 from R2 if the interval could be
asserted without a censoring term --- but that assertion has no ground in a
vocabulary with no term for the flag, which is the distinction at issue rather
than a counter-example to it. The sensitivity is reported in
\texttt{eval/separation.md} rather than left for a reader to raise.

The same three readings run through the reasoner against CENSO entail
\texttt{CensoredObservation} with \texttt{MethodInsufficient},
\texttt{QuantifiedObservation} with \texttt{Exceedance}, and
\texttt{UnresolvedObservation} with no outcome at all. The outcome is entailed
and not asserted: \texttt{Exceedance} is an equivalent class of an observation
bearing \texttt{exceeds} to a threshold. And the separation is a property of
the vocabulary rather than of the loader that populates it --- merging any two
readings of the same individual is an inconsistency, from the disjointness of
the detection statuses. A pipeline cannot quietly hold two readings of one
record.

None of this says the other vocabularies are wrong, or that they could not be
extended; any of them could be given a censoring class tomorrow. It says what
reuse as published currently supports, and where the distinction has to live
once it is added: on the result, not on the device.

\subsection{The ontology answers the questions it was built for}
\label{sec:results-cq}
Twenty competency questions, each tied to an axiom, were executed over the
Waterbase graph with the two vocabulary modules and both regulation packages,
\num{507708} triples as
loaded\src{scripts/17\_run\_competency\_questions.py}. \textbf{All twenty
return answers.} Full OWL~2~RL closure over a graph this size is impractical in
pure Python, so only two entailments are materialised: union membership for the
two defined observation classes, and \texttt{rdfs:subClassOf} on type
assertions. That is measured rather than asserted: over synthetic ABoxes on the
same vocabulary, closure time grows super-linearly in the number of
observations, because RL closure writes every entailment into the graph instead
of answering against rules; SHACL validation over the same graphs costs about an
order of magnitude less\src{scripts/14\_reasoning\_benchmark.py}. The series
is in \texttt{supplementary/S13\_reasoning\_cost.csv}. A deployment over a
whole basin would use a store with native rule support, so those figures bound
this pipeline and not the approach. Every other answer comes from asserted triples. They run on public
data, and their SPARQL, rows and times are in the supplementary material.

\subsection{Four real rows, decided twice}
\label{sec:results-walkthrough}
Figure~\ref{fig:walkthrough} takes one station-year from each of four outcome
classes and decides it twice: once by the vocabulary, once by a two-valued
pipeline at each of the three substitution constants in
use\src{scripts/92\_decision\_walkthrough.py}. The rows are the first match per
class in sorted order, so the selection cannot flatter the argument, and the
two-valued verdicts are computed by the same function that produces the counts
in Section~\ref{sec:results-wbabox} rather than by a second implementation.

\textbf{Two of the four get a different two-valued verdict depending on which
constant the pipeline happened to use}, and nothing in a two-valued output
records that a choice was made. Permethrin is the clearest: a quantification
limit of \num{0.02}\,\si{\micro\gram\per\litre} against a standard of
\num{0.00027} --- \num{74} times the threshold it was meant to decide --- reads
as compliant at zero and as exceeding at half. The first case is included
because both systems agree on it: a bound that clears the standard decides the
question, and a vocabulary that could not say so would be useless.

\begin{figure}[htbp]\centering
\includegraphics[width=\textwidth]{fig11_decision_walkthrough.pdf}
\caption{One real station-year per outcome class, decided by CENSO (left, the
interval the observation establishes against the standard) and by a two-valued
schema (right, at each substitution constant). The shaded bar is what the
measurement supports: for a non-detection it runs from zero, because that is
what ``not detected'' permits. Red line, the standard. Where the right-hand
column changes with the constant, the reported verdict is the convention rather
than the measurement. Every value is read from the shipped graph and ships in
\texttt{supplementary/figure\_data/fig11\_decision\_walkthrough.csv}.}
\label{fig:walkthrough}
\end{figure}

\subsection{The shipped graph is validated against the shipped shapes}
\label{sec:results-shacl}
The published shapes are run over the published graph --- \num{633960} triples
with the vocabulary loaded --- rather than over a fixture\src{scripts/18\_shacl\_validate.py}.
The shapes match \num{40000} assessed observations, \num{31167} censored ones,
\num{1159} analytical methods and \num{537} thresholds, and the run reports
\textbf{one violation type}: \num{4646} observations that cite no analytical
method. That is the unresolvable population, and it is a property of the record
rather than of the pipeline. A graph made to conform would have to invent the
datum, which is the move these shapes were corrected for once already.

Reporting a validation that does not conform requires saying which failures are
findings and which are defects, so the audit is given the list: a violation of a
declared source-data shape is recorded with its count, and any other violation
type fails the build\src{scripts/99\_audit.py}. That distinction had to be made
because until recently this stage ran for over an hour, returned
\texttt{conforms: False}, and was cited by nothing --- not the manuscript, not
the supplementary material, not the audit. The failure it was reporting was
real: the graph's concentration literals were written as bare Turtle
numerals against shapes requiring \texttt{xsd:decimal}, so
\texttt{censo:resultLowerBound 0} carried \texttt{xsd:integer} and could never match the axiom
\texttt{CensoredObservation} $\sqsubseteq$ \texttt{resultLowerBound} value
\texttt{"0.0"\^{}\^{}xsd:decimal} that the vocabulary asserts over every censored
row. No \textsc{owl} check could object, because \texttt{xsd:integer} is derived
from \texttt{xsd:decimal} and satisfies the range.

\subsection{The axioms catch real defects}
Ten test cases run against the vocabulary\src{scripts/test\_axioms.py}. Eight
inject a defect present in this dataset --- a record marked both non-detect and
quantified, a non-detection reported as an exceedance --- and all eight are
rejected, seven by OWL~2~RL entailment and one, outside the profile, by SHACL.
The other two are controls that must remain \emph{consistent}: a well-formed
graph, and one measurement carrying two verdicts under two jurisdictions'
standards. The second control matters because an earlier draft did reject it,
declaring the compliance outcomes disjoint at the class level, which would have
made Section~\ref{sec:results-dual} an inconsistency.

An OOPS!\ scan of the published distribution reports no critical pitfalls.
Three important ones were genuine and are fixed: properties lacking a domain or
range. The rest are artefacts of packaging the modules as one file for the
scanner, which leaves borrowed terms undeclared. Of the thirty-one
properties flagged as lacking an inverse, one genuinely did
(\texttt{definesThreshold}/\texttt{definedBy}) and is now declared; no inverse
was invented for the other thirty.